% ============================================================
%  Poul Martin Møller, "Qvindelighed" (Womanliness)
%
%  Written 1822--1823; published posthumously in
%  Efterladte Skrifter, bd. 5 (C.A. Reitzel, 1855--56), pp. 5--13.
%  Danish source of truth: transcription.tex
%
%  TRANSLATION (English)
% ============================================================

\documentclass[12pt]{article}
\usepackage[utf8]{inputenc}
\usepackage[T1]{fontenc}
\usepackage{libertinus}
\usepackage{libertinust1math}
\usepackage[english]{babel}
\usepackage[protrusion=true,expansion=true]{microtype}
\usepackage[margin=1.4in]{geometry}
\usepackage{setspace}
\usepackage{fancyhdr}
\usepackage{marginnote}
\usepackage[colorlinks=true,linkcolor=black,urlcolor=black,
  pdftitle={Womanliness},
  pdfauthor={Poul Martin Møller}]{hyperref}
\setlength{\headheight}{15pt}
\pagestyle{fancy}
\fancyhf{}
\fancyhead[LE,RO]{\thepage}
\fancyhead[RE]{\textit{Womanliness}}
\fancyhead[LO]{\textit{Poul Martin Møller}}
\setlength{\parindent}{1.5em}
\setlength{\parskip}{0pt}
\onehalfspacing

% Original printed page numbers (1855--56 ed.) as marginal notes
\newcommand{\opage}[1]{\marginnote{\footnotesize\textit{[#1]}}}
% Centred rule marking the author's section breaks
\newcommand{\qvbreak}{\par\medskip\begin{center}\rule{0.32\textwidth}{0.4pt}\end{center}\medskip}

\title{\textbf{Womanliness}}
\author{Poul Martin Møller}
\date{Written 1822--1823; published posthumously in \textit{Efterladte
  Skrifter}, vol.~5,\\ Copenhagen: C.A.\ Reitzel, 1855--56, pp.~5--13}

\begin{document}
\maketitle

\opage{5}To be regarded among women as a model of excellence, one need not so much
perform anything able in the world as give voice to lofty enthusiasm for virtue
and carry fine principles on one's lips. In this respect too, ladies are more
spiritualistic; they make no demand that virtue --- at least not their own ---
should express itself in coarse materialistic deeds. A nobler nature brings them
to keep themselves constantly in the higher region of thought; they love the
moral concepts entirely for their own sake, without requiring them to be stamped
out in the coarse earthly matter. Virtue really loses nothing thereby; for what it
forfeits in intension it regains in extension. The virtuous deeds that a mortal
can perform in his short earthly life amount to a highly meager number. A
gentlewoman, by contrast, can in a single fair evening hour have more rapturous
resolutions and lovely feelings than Aristides expended in his whole lifetime. And
yet virtue is fully present, only as a finer essence, dissolved and \ldots\
[illegible]. My late cousin was a remarkable example of this. Her old aunt, who
spent her last days with her, had been given the household remedy for her
rheumatism that she should be rubbed on the back every evening with a flannel
cloth. She had the peculiarity that no one but my cousin could perform this
service quite to her liking. But unfortunately the latter was wont at just that
same time to be absorbed in tales of noble ladies and nobles who sacrificed their
earthly happiness for the sake of the Kingdom of God. This ravishing reading
\opage{6}she could seldom bring herself to exchange for so clumsy a piece of
handiwork, and the aunt was therefore usually neglected. Once I said to her: Good
cousin! You really ought to take a little more pity on the old soul. In view of
her age she will not long be able to be a trouble to you, and when you have lost
her, you will surely regret your want of attentiveness toward her. --- My friend!
she answered, you do not comprehend how repugnant such petty chores are to me. The
kind of self-conquest that so small a work of love presupposes, I find it hard to
condescend to. I can, as well as anyone, deny myself for my suffering sister and
brother, but only in great things; what anyone can do, I shrink from setting
about. My heart will always be misjudged, because I am born in a wretched age
where I do not belong, in an age when woman is merely to occupy herself with
trifles. Die for my good aunt I can with joy, indeed shed my blood for her on the
scaffold. That I hope you believe of me; but to waste my time humoring her
willfulness, I hold myself too good for.

My cousin really felt herself capable of dying for her aunt, indeed for almost
anyone. But by this I by no means wish to have said that this lofty fancy would
really, given the occasion, have been carried into effect here in this world. Had
a filthy knacker's man, whose boots reeked of train-oil, with unwashed hand and
unclipped nails, really seized her by her fair hair, she would surely have formed
a better resolution in time. It was really the pure representation --- to die for
one's neighbor, the fine concept, stripped of all crude additions --- that
inspired her mind. She would not have remained true to her intention had she seen
the death-scene depicted true to nature in a black copper engraving, still less if
the picture had been hand-colored.

She shunned as much as possible all contact with \opage{7}prosaic reality; for
experience daily made a new conquest from the realm of thought in which she
performed her ideal heroic deeds. In a solitary room, on her soft sofa, when quite
undisturbed she gave herself over to the dream of what she would have done under
other circumstances, she drove her demands upon the human being so high that she
with good reason despised Leonidas, Winkelried, and Jeanne d'Arc; for she knew
that in their place she would have done more than they did. Only for Mucius
Scaevola did she have a secret respect; she felt he was her superior. The reason
for this was that she had once burned herself on a curling iron, and had therefore
solemnly promised herself never willingly to come too near the flame. Thereby her
world of fancy had been restricted to three elements. But was she less high-minded
because her health did not permit her to stand in interaction with the visible? If
moral worth were to be judged by the ability with which one knows how to expose
oneself to the inconvenience of outward things, then surely the fishwife, who sits
from morning to evening in rain and sleet by her basket without getting a
toothache, would have to be given a particularly high place on morality's
rank-list. When each person works in the sphere in which he feels at home, one
cannot demand more of him. One feels able to tumble about with wet cod and
whiting, another with thoughts and representations, according as each one's
constitution allows. But is one not tempted by such reflections to wish that the
Creator had, for the use of finer beings, prepared a finer stage of distilled
elements; that he had provided a more convenient pocket edition of the world for
the use of nervous persons?

I touched earlier on my cousin's demands upon her neighbor's morality. On this I
must declare myself somewhat more definitely. In this respect she was very
reasonable, and did not demand much more of them than of herself. She did not
require of her neighbor that he \opage{8}should be strict in the fulfillment of
any duties other than those he had toward her. For the rest she let herself be
completely satisfied provided he only had a fine theory of life. Noble principles
he had to possess, and have talent for bringing them to light in a fine language;
then he was called, in her terminology: a dear one. Even if one told of such a
person that he was an idler and did not do a penny's worth of good here in the
world, she nonetheless remained firm in her devotion to him, and always gave the
excusing answer: Yes, but he is after all such a dear one; or else: Yes, but he is
after all so sweet.

One of our mutual friends stood in very bad repute with my cousin; for he is a
great man of practice, whose time does not permit him to occupy himself much with
lofty views of life. Since he is reluctant to sketch any plan that he does not see
himself able to carry out, he sets himself firmly against all high-pitched demands
upon weak mortals. The consequence of this was that he would rather curtly dismiss
Adelaide when she sought to draw him into her rapturous circle of thought; indeed,
with a kind of defiance he often gave voice to a far more trivial philosophy of
life than the one he himself follows. This man was naturally most repugnant to
her. I now and then took him into my defense and begged her to note that in
reality he daily showed more self-sacrifice than most. But here she again
displayed her ideal nature, since she set the word and the concept high above
coarse deeds. He was and remained to her utterly insufferable.

\qvbreak

In general, moralists inveigh strictly against women who give up unusually much of
their bodily beauty to the eyes of curious onlookers. That is not very becoming;
but there is another kind of exposure, which is generally praised, and which yet
at bottom conflicts far more with true womanliness. \opage{9}A spirited woman has
seen the matter in its right light when she says that every true feeling is
accompanied by a certain modesty. This belief a good many of her sisters seem not
at all to profess, since they precisely with diligence display their fine feelings
quite plainly and use the exhibition of them as a kind of poetic ornament for
their person. But if it betrays a want of modesty not duly to cover the outer
person, then it certainly presupposes an even greater disregard of modesty
deliberately to bare the inner person; to make a hunt for applause by an
exhibition of the heart's noble character. --- A certain degree of culture makes
the unconscious outbursts of feeling highly rare among human beings, but such
feeling had rather remain hidden in its private chamber than step forth with
self-conscious theatrical phrases and gestures. When a woman knowingly displays
her inner emotions, she necessarily, by her nature, thinks of whether it becomes
her prettily or not. From this it follows that her conduct on such an occasion
becomes a kind of refined coquetry. Secretly every woman herself feels the true
significance of such behavior; for when she has not much practice in giving such
performances, she is apt to blush somewhat at her own bookish pathos. But a truly
innocent maiden almost never lays her emotions bare in company of several; they
are hidden deep in the secrecy of her heart, or poured out to a woman friend or a
mother (not to a sister), and outwardly she is surrounded by an apparent severity
and coldness, which only a fool mistakes for heartlessness.

\qvbreak

\begin{center}\large G.~F.~G.\end{center}

You must not be surprised that I begin at once, without introduction, upon the
text of this letter. My intention is \opage{10}to explain more closely a couple of
words that I, boldly but loosely, let fall yesterday evening in conversation with
you and your sister. It is hardly possible for me to make myself fully intelligible
to you; for as a consequence of the age's fragmentary development, one person in
our days is almost never rightly understood by another. But I feel nonetheless
that I might possibly be much misunderstood were I to let the matter rest with our
oral dispute.

First, then, I must seek to make good that I have by no means wished to come too
near woman's honor by asserting, what I still assert, that it is \emph{unbecoming},
indeed \emph{abominable}, that a woman should be a poetess by \emph{profession}. By
this I set them neither above nor below men, since I consider it most preposterous
to want to apply such an ordinance of precedence to our Lord's garden. At most one
can only compare things of the same kind; more closely regarded, perhaps not even
that. This much then becomes certain, that one cannot determine whether, for
example, an oak or a beech stands highest on our Lord's list, whether a thunderclap
or a poplar is highest, \&c.\ \&c. But that work of nature is nonetheless the most
beautiful which comes nearest to Nature's thoughts, which most purely presents the
Creator's idea. It is therefore incontestably most beautiful that a woman be truly
womanly. When, then, on one's way through life one meets a Miss ***, one has surely
no right to criticize the bizarre fancy to which Nature has now and then abandoned
herself in strange combinations of conflicting elements; but one never calls such a
phenomenon beautiful. Precisely the same effect it has on me to think of a woman as
a scholar or as --- a poetess!! I will not appeal to history, though one learns from
it that all poetesses have been spiritual deformities and freaks of nature (Sappho,
Mad.\ Dacier, Mad.\ de Staël, Mrs.\ B.\ \&c.). But it highly offends my immediate
tact for the womanly-beautiful to see a woman chew \opage{11}on her pen or sit with
her folio in her lap. It puts me in just the same mood as if I saw her play cards
and trump vigorously on the table, swear loud and hard, smoke her cigar, walk about
in waterproof boots, and, with a brisk mannish spitting, describe great arcs in the
air. There is a certain active kind of conduct in life --- to which I reckon
authorship, versifying, preaching, brawling --- which to my mind ill becomes women.
Were you very orthodox, I could appeal to the Apostle Paul, who says that woman
shall keep silent in church.

If we consider the essence of poetry more closely, it becomes even clearer as day.
There are only two kinds of poetry. Either historical memory or the present age can
furnish the material. Surely the highest kind of art is that which arises as an
echo of a glorious vanished age, whose memory lives fresh in beautiful legends on
all lips. Then some particularly highly gifted persons must naturally arise, who
most purely apprehend the image of the age's completed drama. For this that
Providence takes care, which lets nothing magnificent pass away without a memorial.
This kind of art, in which the offspring of history, purified by the age's
purgatorial fire, again step forth with a renewed life, is a pure product of
nature, like the luxuriant growths that swarm forth from fallen ruins. To this
class belong the kings of all poets: Homer, Shakespeare, and Ariosto. Here women
naturally do not come into consideration at all. When has a woman revived the
memory of antiquity? When has she even \emph{attempted} to bring forth anything, in
the proper sense, high or great in poetry?

The other branch of poetry is the personal. It is the only real art there is in our
age. (For the weak, half-whispering echoes from ages that have long since received
their vigorous funeral sermon do not deserve to be counted at all.) It follows of
itself that one whose life has been stirred by strong passions comes, when the
storm has subsided, to composure and can, \opage{12}if he possesses the gift of
presentation, give us an interesting presentation of what he suffered and lived
through. All able poetic works in our age, however disguised they may come forth,
have their root in the authors' own lives. Goethe himself acknowledges that Werther
and Faust are fruits of his own passionate youth, and only thereby do they win that
wonderfully ravishing truth. However great a spirit Goethe is, he has never wished
to live in order to be a poet. His poetic works came into the world under his
manifold studies, almost without decision, involuntarily. Nowadays to live in order
to be a poet is the same as to travel in order to produce a travel-book. It always
looks ill to feel in order to depict one's feelings; but doubly so when ladies do
it. They cannot, as a consequence of life's present tight forms, even come to lead
so motley a life as one must have led who is to have anything remarkable to depict.
And those who have done so are ladies of a corresponding stamp. The most innocent
and most common ladies' poetry consists in imitation of what others have created.
But such products signify, in the realm of art, no more than the ordinary
black-chalk copies of Petrarch and Laura, Achilles, Homer, Ossian, \&c., which one
sees hanging by a green silk ribbon on the parlor walls. But among those women in
whom poetry really is something --- that is, in whom it is spun out of their own
breast --- it is far more \emph{culpable}. Such a poetess must anxiously lie in wait
for all her feelings, in order at once to get them entered in the books. None of
her joys, none of her tears goes to waste; everything is neatly registered in rhyme
and put on display for the public. When she embraces her mother, she thinks of the
beautiful attitude she thereby gains for her poetic lumber-room; when she strolls
out into God's free nature, she has, ever so cunningly, a bit of vellum paper in her
sewing-bag, so as to have the tools at hand at once, should she be able to work up
as much emotion as goes \opage{13}into a Matthisson-like poem. Should her lover
become faithless (which neither God nor any man, I dare say, blames him for), then
she rejoices secretly that she nonetheless thereby profits a score of sonnets, or
some essential ingredients for a doleful novella. In short, all her life through she
sits before her poetic mirror and makes silhouettes of her little feelings the
moment they have half come into the world.

This, good madam, is my humble opinion of poetesses, and in accordance with it, it
will not surprise you that I rejoice that we do not have even more of that sort of
people in our country. Now I will beg you not to become angry that I have laid claim
to your time with these lines. I should not have ventured upon it, had I not, from
your judgment on sermons, been brought to the opinion that you are not so easily
bored. Seriously speaking, I should like to efface one of the disagreeable
impressions that I sometimes cannot help making upon you by the hard, coarse,
clumsy manner in which I so often, against my will, come to express myself in your
presence. I really admire the infinite gentleness and good-naturedness with which
you put up with it. For some time now I have, quite against my will, acquired that
sharp manner in life which often afterward sets even myself in astonishment.

\begin{flushright}
Your devoted\\
P.~Møller.
\end{flushright}

\qvbreak

\end{document}
